\textbf{Время:} резерв

Использование: дополнительный материал для вопросов после основной части. Температура и химический потенциал не нужны для понимания основной цепочки измерения.

\textbf{Тезис:} низкоэнергетический режим STAR FXT исследует барионно насыщенную материю и дополняет картину КГП при более высоких энергиях.

\textbf{Визуал:} вертикальная ось — температура T, горизонтальная — барионный химический потенциал μB. КГП подписана в верхней части, адронная материя — в нижней. Полоса справа показывает интерес к большим μB, а не измеренную область температур при 3 ГэВ. Пунктир обозначает неопределённость продолжения линии перехода. На схеме нет численной траектории и установленной критической точки.

\textbf{Текст:} При высокой температуре обычные адроны уже не являются удобным описанием всех степеней свободы: говорят о кварк-глюонной плазме. Но у фазовой карты есть и второе направление. Химический потенциал связан с барионным избытком среды; это не простое другое название плотности. При понижении энергии столкновения сильнее проявляется остановка исходных барионов. Три ГэВ позволяют исследовать именно эту барионно насыщенную область.

Это не означает, что в столкновениях при 3 ГэВ уже установлено образование КГП. Измерения потока STAR при этой энергии указывают на существенную роль барионных взаимодействий и отсутствие наблюдаемого при более высоких энергиях масштабирования по числу составляющих кварков. Конкретное фазовое состояние требует совокупности наблюдаемых и модельного описания.

Статья Краевой приводит для диапазона FXT 3–7.7 ГэВ ориентир μB около 720–420 МэВ. Это контекст термодинамического описания, а не результат HBT и не координаты каждой стадии столкновения. Фемтоскопия измеряет корреляции конечных пионов и характеризует позднее испускание. Она даёт проверку сценария расширения, а не прямой снимок ранней КГП.

\textbf{Источники:} статья Анны Краевой, Universe 10, 188 (2024), введение и раздел 2; DOI: \url{https://doi.org/10.3390/universe10040188}.

STAR, Disappearance of partonic collectivity in sqrt(sNN)=3 GeV Au+Au collisions at RHIC, Phys. Lett. B 827 (2022) 137003, arXiv:2108.00908, \url{https://arxiv.org/abs/2108.00908}. Схема авторская, качественная; высокая μB не изображена как установленная граница или доказательство КГП.

Для обсуждения: фемтоскопия дополняет другие наблюдаемые, но сама по себе не определяет фазовое состояние ранней материи.
